\chapterbackmatter{Weinberg Exercises}

\begin{enumerate}
\WeinbergExercise{1} Consider a theory with a separable interaction; that is,
  \[
    \langle\beta|H_{\rm int}|\alpha\rangle=g\,u_\beta u_\alpha^*,
  \]
  where $g$ is a real coupling constant, and $u_\alpha$ is a set of complex
  quantities with
  \[
    \sum_\alpha|u_\alpha|^2=1.
  \]
  Use the Lippmann--Schwinger equation \eqref{eq:3.1.16} to find explicit solutions for
  the `in' and `out' states and the $S$-matrix.

\WeinbergExercise{2} Suppose that a resonance of spin one is discovered in $e^+e^-$
  scattering at a total energy of 150 GeV and with a cross-section (in the
  center-of-mass frame, averaged over initial spins, and summed over final
  spins) for elastic $e^+e^-$ scattering at the peak of the resonance equal to
  $10^{-34}\,\mathrm{cm}^2$. What is the branching ratio for the decay of the
  resonant state $R$ by the mode $R\to e^-+e^+$? What is the total
  cross-section for $e^+e^-$ scattering at the peak of the resonance? (In
  answering both questions, ignore the non-resonant background scattering.)

\WeinbergExercise{3} Express the differential cross-section for two-body scattering in the
  \emph{laboratory} frame, in which one of the two particles is initially at
  rest, in terms of kinematic variables and the matrix element
  $M_{\beta\alpha}$. (Derive the result directly, without using the results
  derived in this chapter for the differential cross-section in the
  center-of-mass frame.)

\WeinbergExercise{4} Derive the perturbation expansion \eqref{eq:3.5.8} directly from the expansion
  \eqref{eq:3.5.3} of old-fashioned perturbation theory.

\WeinbergExercise{5} We can define `standing wave' states $\ket{\alpha}^{0}$ by a modified
  version of the Lippmann--Schwinger equation
  \[
    \ket{\alpha}^{0}=\ket{\alpha}_0
    +\operatorname{P.V.}\frac{1}{E_\alpha-H_0}
      H_{\rm int}\ket{\alpha}^{0}.
  \]
  Show that the matrix
  $K_{\beta\alpha}=\pi\delta(E_\beta-E_\alpha)
  \langle\beta|H_{\rm int}|\alpha\rangle^0$ is Hermitian. Show how to express
  the $S$-matrix in terms of the $K$-matrix.

\WeinbergExercise{6} Express the differential cross-section for elastic $\pi^+$--proton and
  $\pi^-$--proton scattering in terms of the phase shifts for states of
  definite total angular momentum, parity, and isospin.

\WeinbergExercise{7} Show that the states $\ket{E,\mathbf p,j,\sigma,\ell,s,n}$ defined by
  Eq. \eqref{eq:3.7.5} are correctly normalized to have the scalar products \eqref{eq:3.7.6}.
\end{enumerate}
